% Reciprocal proofs for the independent formal reconstructions in Goodman projectibility.
\begin{annexitem}{annex:p15-goodman-grue-models}{thm:p15-goodman-grue-models}{Finite grue countermodels}
\begin{proof}
Take $U=D\cup\{b\}$ and interpret $\mathsf{E}$ by $U$ and
$\mathsf{O}_t$ by $D$ in both structures. In $w_G$, interpret
$\mathsf{G}$ by $U$ and $\mathsf{B}$ by the empty set. For every
$a\in D$, the first branch of the grue definition holds, so
$\mathsf{Q}_t(a)$ is true. At $b$ the first branch is false because
$\mathsf{O}_t(b)$ is false; the second is false because
$\mathsf{B}(b)$ is false. Thus the grue extension is exactly $D$.
All emeralds are green but $b$ is an emerald which is not grue.

In $w_Q$, interpret $\mathsf{G}$ by $D$ and $\mathsf{B}$ by
$\{b\}$. Every $a\in D$ is still an examined green emerald and
grue by the first branch. The object $b$ is grue by the second branch.
The grue extension is therefore $U$, while $b$ is not green. All
emeralds are grue but not all are green. The structures agree on every
recorded predicate fact about the objects in $D$, including their
blueness being false. Green and blue are disjoint in both structures.

For the general incompatibility claim, suppose $H_G,H_Q$ both hold
and choose an emerald $b$ with $\neg\mathsf{O}_t(b)$. The first
hypothesis gives $\mathsf{G}(b)$. The second gives
$\mathsf{Q}_t(b)$, whose definition, together with
$\neg\mathsf{O}_t(b)$, gives $\mathsf{B}(b)$. This contradicts the
assumption that an emerald cannot be both green and blue. The witness
$b$ and the incompatibility premise were both used; without them
the incompatibility conclusion is not established.
\end{proof}
\end{annexitem}

\begin{annexitem}{annex:p15-goodman-vocabulary}{thm:p15-goodman-vocabulary}{Inverse Boolean descriptions and name insertion}
\begin{proof}
Fix an arbitrary $x$. If $\mathsf{O}_t(x)$ holds, the grue definition
reduces to $\mathsf{Q}_t(x)\Longleftrightarrow\mathsf{G}(x)$,
and the bleen definition to
$\mathsf{R}_t(x)\Longleftrightarrow\mathsf{B}(x)$.
Consequently the displayed right-hand side for green reduces to
$\mathsf{Q}_t(x)$ and hence to $\mathsf{G}(x)$; the right-hand
side for blue reduces to $\mathsf{R}_t(x)$ and hence to
$\mathsf{B}(x)$.

If $\neg\mathsf{O}_t(x)$ holds, the definitions reduce instead to
$\mathsf{Q}_t(x)\Longleftrightarrow\mathsf{B}(x)$ and
$\mathsf{R}_t(x)\Longleftrightarrow\mathsf{G}(x)$.
The right-hand side for green now reduces to $\mathsf{R}_t(x)$,
and the one for blue to $\mathsf{Q}_t(x)$, again giving the asserted
identities. Classical excluded middle covers all $x$. Neither case
uses a relation between the two colour predicates, so no disjointness
or exhaustiveness condition has been introduced.

In any ordinary first-order structure with equality, $c=c$ is true.
Hence $H\land(c=c)$ is true exactly when $H$ is true. Introducing
$c$ into a language with a nonempty domain permits every expansion
assigning $c$ an arbitrary element; the sentence's truth is independent
of that assignment. Thus its overt use of an individual name does not
distinguish its content from that of $H$.
\end{proof}
\end{annexitem}

\begin{annexitem}{annex:p15-goodman-confirmation-collapse}{thm:p15-goodman-confirmation-collapse}{Collapse of the unrestricted confirmation relation}
\begin{proof}
Choose arbitrary formulas $e,h$. The conjunction $e\land h$ entails
$e$, so the first assumed condition gives $C(e,e\land h)$.
The same conjunction entails $h$, so the second assumed condition
gives $C(e,h)$. No restriction has been placed on either chosen
formula. Therefore the relation contains every ordered pair.
This includes inconsistent formulas if they belong to its domain;
no such formula was needed to obtain the result. Restricting either
closure condition can invalidate the argument at its identified use,
so the theorem asserts no collapse for a relation lacking one of
the unrestricted premises.
\end{proof}
\end{annexitem}
